\documentclass[12pt]{exam}
\usepackage{amsthm}
\usepackage{bm}
\usepackage{libertine}
\usepackage[utf8]{inputenc}
\usepackage[margin=0.5in]{geometry}
\usepackage{amsmath,amssymb}
\usepackage{multicol}
\usepackage[shortlabels]{enumitem}
\usepackage{siunitx}
\usepackage{physics}
\usepackage{booktabs}
\usepackage{graphicx}
\usepackage{pgfplots}
\usepackage{listings}
\usepackage{tikz}
\usepackage{tikz-3dplot}

\bmdefine{\ii}{i}
\bmdefine{\jj}{j}
\bmdefine{\kk}{k}
\newcommand{\te}{\theta}
\newcommand{\cC}{\mathcal{C}}
\newcommand{\word}[1]{\quad\text{#1}\quad}

\pgfplotsset{width=10cm,compat=1.9}
\usepgfplotslibrary{external}
\tikzexternalize

\newcommand{\class}{Math 261-001}
\newcommand{\examnum}{Quiz 1 Solutions}
\newcommand{\examdate}{August 28}

\begin{document}
\pagestyle{plain}
\thispagestyle{empty}

\noindent
\textbf{\class}\\
\textbf{\examnum}, \textbf{\examdate}

\vspace{10pt}
\textbf{Total: 40 points}

\begin{enumerate}

\item \textbf{Expected final answer:}
\[
  2\ii-\jj+3\kk.
\]
\textbf{Verification:} The coefficients of $\ii$, $\jj$, and $\kk$ are
$2$, $-1$, and $3$, respectively.

\textbf{Equivalent forms:} $2\ii+(-1)\jj+3\kk$ and algebraically identical
standard-basis notation receive full credit.

\textbf{Grading (10 points):} 3 points for each correct coefficient and sign;
1 point for valid standard-basis notation.

\item \textbf{Expected final answer:}
\[
  -1.
\]
\textbf{Verification:}
\[
  \vec{u}\cdot\vec{v}=1(3)+2(-1)+(-1)(2)=3-2-2=-1.
\]
\textbf{Equivalent forms:} Any numerical expression equal to $-1$ receives
full credit.

\textbf{Grading (10 points):} 5 points for the three component products,
3 points for combining them correctly, and 2 points for the final value.

\item \textbf{Expected final answer:}
\[
  \vec{r}(t)=(1,-1,2)
  +t(2,2,-1),\qquad t\in\mathbb{R}.
\]
\textbf{Verification:} $Q-P=(2,2,-1)$; the formula gives $P$
at $t=0$ and $Q$ at $t=1$.

\textbf{Equivalent forms:} Any affine parametrization
$\vec{r}(t)=t\vec{v}+\tilde{P}$ of the same line receives full credit: $\tilde{P}$ must
lie on the displayed line and $\vec{v}$ must be a nonzero scalar multiple
of $(2,2,-1)$. Equivalent component equations are also valid.

\textbf{Grading (10 points):} 4 points for a valid direction vector,
4 points for a point on the line, and 2 points for a complete parametric
equation.

\item
\begin{enumerate}[(a)]
  \item \textbf{Expected final answer:}
  \[
    \overrightarrow{PQ}\times\overrightarrow{PR}=(1,-1,1).
  \]
  \textbf{Verification:}
  \[
    \overrightarrow{PQ}=(1,1,0),\qquad
    \overrightarrow{PR}=(0,1,1),
  \]
  \[
    \overrightarrow{PQ}\times\overrightarrow{PR}
      =(1(1)-0(1),\;0(0)-1(1),\;1(1)-1(0))
      =(1,-1,1).
  \]
  \textbf{Equivalent forms:} Ordered-triple and standard-basis notation are
  both valid. Because the order of the cross product is specified, a different
  nonzero scalar multiple is not the requested cross product.

  \textbf{Grading (5 points):} 2 points for the two edge vectors and 3 points
  for the correct cross product.

  \item \textbf{Expected final answer:}
  \[
    x-y+z=2.
  \]
  \textbf{Verification:} With $\vec{n}=(1,-1,1)$ and $P$,
  \[
    \vec{n}\cdot(x-1,y,z-1)=0
    \quad\Longleftrightarrow\quad x-y+z=2.
  \]
  Substitution confirms that $P$, $Q$, and $R$ each satisfy the equation.

  \textbf{Equivalent forms:} Any algebraically equivalent equation, including
  any nonzero scalar multiple of $x-y+z-2=0$, receives full credit.

  \textbf{Grading (5 points):} 2 points for using a valid normal, 1 point
  for substituting a point correctly, and 2 points for a correct plane
  equation.
\end{enumerate}

\end{enumerate}

\end{document}
